% --- Archivo: 01_subespacio_tangente.tex ---

Este capítulo presenta las condiciones teóricas de optimalidad para problemas con restricciones de igualdad, definidos sobre variedades diferenciables. A partir del Teorema de Lyusternik, se establece la relación entre la topología del cono tangente y el núcleo de la matriz Jacobiana, para luego deducir las condiciones de primer y segundo orden mediante multiplicadores de Lagrange.

\begin{notation}
    Dada una función vectorial diferenciable $h : \Rn \to \R^l$, denotaremos por $J_h(x) \in \R^{l \times n}$ a su matriz Jacobiana evaluada en el punto $x$. Consecuentemente, las filas de esta matriz corresponden a los gradientes transpuestos $\nabla h_i(x)^\top$, para $i = 1, \dots, l$.
\end{notation}

% =========================================================
\section{El Problema de Optimización y el Subespacio Tangente}
% =========================================================

\begin{definition}[Problema con Restricciones de Igualdad]\label{def:problema_igualdad}
    Sean $f : \Rn \to \R$ y $h : \Rn \to \R^l$ funciones de clase $\mathcal{C}^1$ (y de clase $\mathcal{C}^2$ cuando el análisis lo requiera). Se define el \textbf{conjunto factible} $D$ como la variedad algebraica descrita por el sistema de igualdades:
    \[
        D = \{ x \in \Rn \mid h(x) = 0 \}.
    \]
    El \textbf{problema de optimización con restricciones de igualdad} consiste en hallar un elemento en $D$ que minimice a la función objetivo $f$. Este problema se denota formalmente mediante la expresión:
    \begin{equation}
        \min f(x) \quad \text{sujeto a} \quad x \in D.
        \label{eq:problema_igualdad}
    \end{equation}
\end{definition}

\begin{proposition}\label{prop:bouligand_en_kernel}
	Supóngase que $h : \Rn \to \R^l$ es diferenciable en $\bar{x} \in D$. Entonces, el cono de Bouligand está contenido en el núcleo del Jacobiano:
	\[
		\mathcal{B}_D(\bar{x}) \subset \ker J_h(\bar{x}).
	\]
\end{proposition}

\begin{proof}
	El vector nulo pertenece a ambos conjuntos. Sea $d \in \mathcal{B}_D(\bar{x}) \setminus \{0\}$. Por definición secuencial, existen sucesiones $\{t_k\} \to 0^+$ y $\{d^k\} \to d$, tales que $\bar{x} + t_k d^k \in D$ para todo $k \in \N$. Consecuentemente, $h(\bar{x} + t_k d^k) = 0$. Aplicando el desarrollo de Taylor de primer orden a la función vectorial $h$:
	\[
        0 = h(\bar{x} + t_k d^k) - h(\bar{x}) = t_k J_h(\bar{x})d^k + o(t_k \norm{d^k}).
	\]
	Dividiendo por el escalar $t_k > 0$ y tomando el límite cuando $k \to \infty$, la convergencia $d^k \to d$ implica que $J_h(\bar{x})d = 0$. Por consiguiente, $d \in \ker J_h(\bar{x})$.
\end{proof}

\begin{notabox}[Interpretación Geométrica del Núcleo]
    El álgebra lineal indica que $d \in \ker J_h(\bar{x})$ equivale a $\interno{\nabla h_i(\bar{x})}{d} = 0$ para cada restricción $i$. Geométricamente, esto significa que cualquier movimiento en la dirección de $d$ debe ser ortogonal a los gradientes de las superficies que definen a $D$. De este modo, $d$ es tangente a la superficie en primer orden.
\end{notabox}

\begin{definition}[Condiciones de Regularidad]\label{def:condiciones_regularidad}
	Decimos que el punto $\bar{x} \in D$ es \textbf{regular} (o satisface una calificación de restricciones) si cumple una de las siguientes propiedades:
	\begin{enuthm}
		\item \textbf{Independencia Lineal de los Gradientes (LICQ):} El operador lineal $J_h(\bar{x})$ es sobreyectivo; esto es, $\rk J_h(\bar{x}) = l$. Algebraicamente, esto equivale a afirmar que la transformación transpuesta es inyectiva: $\ker(J_h(\bar{x})^\top) = \{0\}$. (Nótese que esto exige la condición topológica $l \le n$).
		\item \textbf{Linealidad afín:} La función $h$ es afín, i.e., $h(x) = Ax - a$, con $A \in \R^{l \times n}$ y $a \in \R^l$.
	\end{enuthm}
\end{definition}

Para establecer la topología local de la variedad, se requiere un resultado analítico preparatorio basado en el Teorema de la Función Implícita clásico.

\begin{lemma}[Existencia de Función Implícita Acotada]\label{lem:funcion_implicita_acotada}
    Sean $F : \R^s \times \Rn \to \R^l$ diferenciable en una vecindad de $(\bar{p}, \bar{y})$ y $\nabla_y F$ continua en dicho punto. Supóngase que $F(\bar{p}, \bar{y}) = 0$ y que el rango de la matriz Jacobiana parcial es máximo: $\rk J_y F(\bar{p}, \bar{y}) = l$.
    
    Entonces, existen una vecindad $U$ de $\bar{p}$ y una constante $M > 0$ tales que, para todo $p \in U$, existe un único $y(p) \in \Rn$ que satisface $F(p, y(p)) = 0$ y la acotación:
    \[
        \norm{y(p) - \bar{y}} \le M \norm{F(p, \bar{y})}.
    \]
\end{lemma}

\begin{proof}
    Como $J_y F(\bar{p}, \bar{y}) \in \R^{l \times n}$ tiene rango $l$, sus filas son linealmente independientes. Sea $A \in \R^{(n-l) \times n}$ una matriz cuyas filas completan una base de $\Rn$. Construimos el operador extendido $\Phi : \R^s \times \Rn \to \R^l \times \R^{n-l}$ mediante:
    \[
        \Phi(p, y) = \begin{pmatrix} F(p, y) \\ A(y - \bar{y}) \end{pmatrix}.
    \]
    Evaluando en el punto base, $\Phi(\bar{p}, \bar{y}) = 0$. El Jacobiano parcial respecto a $y$ es la matriz cuadrada $J_y \Phi(\bar{p}, \bar{y}) = \begin{pmatrix} J_y F(\bar{p}, \bar{y}) \\ A \end{pmatrix} \in \R^{n \times n}$, la cual es no singular por construcción.
    
    Por el Teorema de la Función Implícita, existe una función diferenciable $y : U \to \Rn$ definida en una vecindad de $\bar{p}$ tal que $y(\bar{p}) = \bar{y}$ y $\Phi(p, y(p)) = 0$. Esto último implica que $F(p, y(p)) = 0$ y $A(y(p) - \bar{y}) = 0$.
    
    Aplicando el Teorema del Valor Medio sobre el operador inverso local $\Phi(p, \cdot)$ y utilizando la desigualdad triangular para el difeomorfismo continuo, se deduce la existencia de una constante $M$ dependiente de $\norm{J_y \Phi^{-1}}$ que satisface la acotación requerida.
\end{proof}

\begin{notabox}[Justificación Geométrica de la Corrección Ortogonal]
    Si un punto $x$ no cumple las restricciones ($h(x) \neq 0$), se busca una corrección $\Delta x$ tal que $h(x + \Delta x) = 0$. Geométricamente, la trayectoria más corta para retornar a la variedad es ortogonal a ella. Dado que el espacio ortogonal al plano tangente ($\ker J_h$) coincide con la imagen de la matriz transpuesta ($\im J_h^\top$), la corrección óptima tiene la forma $\Delta x = J_h(\bar{x})^\top \xi$, donde $\xi \in \R^l$ es un vector de ponderación. Este concepto fundamenta el desarrollo del Teorema de Lyusternik.
\end{notabox}

\begin{example}\label{ex:funcion_implicita_acotada}
    Para ilustrar el Lema de la Función Implícita Acotada, consideremos la intersección de una circunferencia de radio paramétrico $p$ con la recta vertical $y_1 = 1$. Modelamos esto mediante el campo vectorial $F : \R \times \R^2 \to \R^2$ dado por:
    \[
        F(p, y) = \begin{pmatrix} y_1^2 + y_2^2 - p^2 \\ y_1 - 1 \end{pmatrix}.
    \]
    Fijemos el parámetro nominal $\bar{p} = \sqrt{2}$ y el punto de solución nominal $\bar{y} = (1, 1)^\top$. Se verifica de manera directa que $F(\sqrt{2}, (1,1)) = 0$.
    
    Evaluamos el Jacobiano parcial respecto a las variables espaciales $y$:
    \[
        J_y F(p, y) = \begin{pmatrix} 2y_1 & 2y_2 \\ 1 & 0 \end{pmatrix} \implies J_y F(\bar{p}, \bar{y})
        = \begin{pmatrix} 2 & 2 \\ 1 & 0 \end{pmatrix}.
    \]
    El determinante de esta matriz es $-2 \neq 0$. Por lo tanto, $\rk J_y F(\bar{p}, \bar{y}) = 2 = l$, por lo que las hipótesis del \cref{lem:funcion_implicita_acotada} se cumplen.
    
    Si se perturba el radio de la circunferencia ligeramente, pasando a un nuevo parámetro $p \approx \sqrt{2}$ y manteniendo el punto en $\bar{y} = (1,1)$, el error algebraico inducido es:
    \[
        F(p, \bar{y}) = \begin{pmatrix} 1^2 + 1^2 - p^2 \\ 1 - 1 \end{pmatrix} = \begin{pmatrix} 2 - p^2 \\ 0 \end{pmatrix}.
    \]
    El lema asegura que existe una raíz exacta $y(p)$ para el sistema perturbado, y que la distancia desde el punto antiguo $\bar{y}$ a la nueva raíz $y(p)$ está acotada por este error: $\norm{y(p) - \bar{y}} \le M \norm{F(p, \bar{y})} = M \absolute{2 - p^2}$. En este caso, la nueva raíz analítica es $y(p) = (1, \sqrt{p^2 - 1})^\top$, y mediante el desarrollo de Taylor se verifica que la distancia $\absolute{\sqrt{p^2 - 1} - 1}$ decae a la misma velocidad que la perturbación $\absolute{p^2 - 2}$.
\end{example}

\begin{theorem}[Lyusternik: Estimación de Viabilidad]\label{thm:lyusternik}
	Supóngase que $h : \Rn \to \R^l$ es continuamente diferenciable en una vecindad de $\bar{x} \in D$ y que satisface la condición LICQ. Entonces, existen una vecindad $U$ de $\bar{x}$ y una constante $M > 0$ tales que:
	\[
		\dist(x, D) \le M \norm{h(x)} \quad \forall x \in U.
	\]
\end{theorem}

\begin{proof}
    Definimos el campo paramétrico $F : \Rn \times \Rn \to \R^l$ mediante la traslación $F(x, \xi) = h(x + \xi)$. Tratando a $x$ como el parámetro y a $\xi$ como la variable, evaluamos en el punto base $(\bar{x}, 0)$. Se tiene $F(\bar{x}, 0) = h(\bar{x}) = 0$. El Jacobiano respecto a la variable $\xi$ es $J_\xi F(\bar{x}, 0) = J_h(\bar{x})$, el cual, por la hipótesis LICQ, tiene rango máximo $l$. Las hipótesis del \cref{lem:funcion_implicita_acotada} se verifican.
    
    Por consiguiente, existen una vecindad $U$ de $\bar{x}$ y un número $M > 0$ tales que, para cada $x \in U$, existe una perturbación $\xi(x) \in \Rn$ que cumple $0 = F(x, \xi(x)) = h(x + \xi(x))$. Esto asegura que el punto trasladado $x + \xi(x) \in D$. Además, el \cref{lem:funcion_implicita_acotada} establece la cota $\norm{\xi(x)} \le M \norm{F(x, 0)} = M \norm{h(x)}$.
    Por la definición métrica de distancia a un conjunto:
    \[
        \dist(x, D) \le \norm{(x + \xi(x)) - x} = \norm{\xi(x)} \le M \norm{h(x)}.
    \]
\end{proof}

\begin{example}[Acotación de Lyusternik en la Esfera]\label{ex:lyusternik_esfera}
    Considérese la variedad esférica unitaria en $\R^2$, definida por $h(x) = x_1^2 + x_2^2 - 1 = 0 \implies D = \{ x \in \R^2 \mid \norm{x}^2 = 1 \}$.
    Tomemos un punto nominal en el polo norte, $\bar{x} = (0,1) \in D$. El Jacobiano evaluado en este punto es $J_h(\bar{x}) = (0, 2)$. Dado que $\rk J_h(\bar{x}) = 1$, la condición LICQ se cumple en $\bar{x}$. El \cref{thm:lyusternik} asegura la existencia de una constante $M > 0$ y una vecindad $U$ de $\bar{x}$ tal que $\dist(x, D) \le M \absolute{h(x)}$ para todo $x \in U$.
    
    Analíticamente, para $x \in \R^2 \setminus \{0\}$, la distancia es $\dist(x, D) = \absolute{\norm{x} - 1}$. La violación algebraica es $\absolute{h(x)} = \absolute{\norm{x}^2 - 1} = \dist(x, D) \cdot (\norm{x} + 1)$. Despejando, obtenemos $\dist(x, D) = \frac{1}{\norm{x} + 1} \absolute{h(x)}$. En una vecindad $U$ alrededor de $\bar{x} = (0,1)$, la norma satisface $\norm{x} \approx 1$. Si restringimos $U$ a la bola de radio $1/2$ centrada en $\bar{x}$, para todo $x \in U$ se cumple que $\norm{x} \ge 1/2$, lo que implica $\frac{1}{\norm{x} + 1} \le \frac{2}{3}$.
    Así, se obtiene la constante de Lyusternik $M = 2/3$. En este caso, el error algebraico $\absolute{h(x)}$ acota la distancia geométrica al conjunto.
\end{example}

\begin{warningbox}[Colapso de la Geometría por falta de Regularidad]
    La condición LICQ es un requisito fundamental para asegurar la regularidad del punto. Consideremos la restricción escalar $h(x) = x^2 = 0$, con el conjunto $D = \{0\}$. En $\bar{x} = 0$, el gradiente se anula ($h'(0) = 0$), por lo que no se cumple LICQ. Para cualquier $x \notin D$, se tiene:
    \[
        \frac{\dist(x, D)}{\norm{h(x)}} = \frac{\absolute{x}}{x^2} = \frac{1}{\absolute{x}} \to +\infty \quad \text{cuando } x \to 0.
    \]
    No existe ninguna constante $M$ finita que acote esta relación, por lo que la equivalencia entre la caracterización algebraica y la topológica deja de cumplirse.
\end{warningbox}

\begin{theorem}[Teorema del Subespacio Tangente]\label{thm:subespacio_tangente}
	Sea $h : \Rn \to \R^l$ de clase $\mathcal{C}^1$ en $\bar{x} \in D$.
	\begin{enuthm}
		\item Si se satisface la condición LICQ en $\bar{x}$, entonces los conos topológicos y el subespacio algebraico coinciden:
		\[
			\mathcal{T}_D(\bar{x}) = \mathcal{B}_D(\bar{x}) = \ker J_h(\bar{x}).
		\]
		\item Si las restricciones son afines, se tiene adicionalmente que $\mathcal{V}_D(\bar{x}) = \ker J_h(\bar{x})$.
	\end{enuthm}
\end{theorem}

\begin{proof}
    \textbf{Parte (1):} Por la \cref{prop:bouligand_en_kernel}, $\mathcal{B}_D(\bar{x}) \subset \ker J_h(\bar{x})$. Inversamente, sea $d \in \ker J_h(\bar{x})$. Por el Teorema de Lyusternik (\cref{thm:lyusternik}), para $t > 0$ suficientemente pequeño se tiene $\dist(\bar{x} + td, D) \le M \norm{h(\bar{x} + td)}$. Expandiendo $h$ por Taylor y utilizando que $h(\bar{x}) = 0$ y $J_h(\bar{x})d = 0$:
	\[
		h(\bar{x} + td) = h(\bar{x}) + t J_h(\bar{x})d + o(t) = o(t).
	\]
    Consecuentemente, $\dist(\bar{x} + td, D) \le M \norm{o(t)} = o(t)$. Esto corresponde a la definición de dirección tangente, por lo que $d \in \mathcal{T}_D(\bar{x})$. De la cadena de inclusiones $\ker J_h(\bar{x}) \subset \mathcal{T}_D(\bar{x}) \subset \mathcal{B}_D(\bar{x}) \subset \ker J_h(\bar{x})$, se deduce la igualdad de los conjuntos.

    \textbf{Parte (2):} Para $h(x) = Ax - a$, se tiene $J_h(\bar{x}) = A$. Si $d \in \ker A$, para cualquier escalar $t \in \R$: $h(\bar{x} + td) = A(\bar{x} + td) - a = (A\bar{x} - a) + tAd = h(\bar{x}) = 0$. Luego $\bar{x} + td \in D$ para todo $t$, lo que demuestra que $d \in \mathcal{V}_D(\bar{x})$.
\end{proof}

\begin{example}[Colapso Analítico del Cono Tangente en $\R^3$]\label{ex:subespacio_tangente}
    Considérese la curva en $\R^3$ formada por la intersección de un cilindro parabólico y un plano, definida por el sistema algebraico $h(x) = (x_1^2 - x_2, x_1 + x_3 - 1)^\top = 0$. Elegimos el punto $\bar{x} = (1, 1, 0)^\top \in D$. El subespacio nulo del Jacobiano evaluado en $\bar{x}$ está generado por el vector director $d = (1, 2, -1)^\top$. Es decir: $\ker J_h(\bar{x}) = \Span{ (1, 2, -1)^\top}$.
    
    Verificamos que este vector pertenece a $\mathcal{T}_D(\bar{x})$. Construimos una curva $x(t) \in D$ parametrizada que pase por $\bar{x}$. Tomando $x_1 = 1 + t$, el sistema exige $x_2(t) = 1 + 2t + t^2$ y $x_3(t) = -t$. La curva factible es $x(t) = (1+t, \, 1+2t+t^2, \, -t)^\top$. Desarrollando: $x(t) = \bar{x} + td + r(t)$, donde el residuo es $r(t) = (0, t^2, 0)^\top$. La norma del residuo satisface $\norm{r(t)} = o(t)$. Por definición, se tiene que $d \in \mathcal{T}_D(\bar{x})$.
\end{example}

\begin{theorem}[Condición Necesaria en Forma Primal]\label{thm:cn_primal_igualdad}
    Sea $\bar{x} \in D$ un minimizador local del problema \eqref{eq:problema_igualdad}. Si $\bar{x}$ es regular (cumple LICQ o linealidad), entonces el gradiente de la función objetivo se anula a lo largo de todo el subespacio tangente:
    \begin{equation}
        \interno{\nabla f(\bar{x})}{d} = 0 \quad \forall d \in \ker J_h(\bar{x}).
    \end{equation}
\end{theorem}

\begin{proof}
    De acuerdo con la caracterización general del \cref{chap:intro_opt}, si $\bar{x}$ es un mínimo local, se cumple que $\interno{\nabla f(\bar{x})}{d} \ge 0$ para todo $d \in \mathcal{B}_D(\bar{x})$. Por la regularidad y el \cref{thm:subespacio_tangente}, se tiene $\mathcal{B}_D(\bar{x}) = \ker J_h(\bar{x})$. Esta desigualdad se aplica sobre el subespacio nulo. Dado que $\ker J_h(\bar{x})$ es un subespacio vectorial, si $d \in \ker J_h(\bar{x})$, su opuesto simétrico $-d$ también pertenece a $\ker J_h(\bar{x})$. Evaluando la desigualdad en dicho vector: $\interno{\nabla f(\bar{x})}{-d} \ge 0 \implies \interno{\nabla f(\bar{x})}{d} \le 0$. De ambas desigualdades se deduce la igualdad a cero.
\end{proof}

\begin{example}[Ortogonalidad Primal en el Cilindro]\label{ex:cn_primal}
    Considérese el problema de minimizar la distancia al origen de un punto restringido a yacer sobre un cilindro desplazado en $\R^3$: $\min \frac{1}{2}(x_1^2 + x_2^2 + x_3^2)$ sujeto a $h(x) = (x_1 - 2)^2 + x_2^2 - 1 = 0$. Geométricamente, el punto del cilindro más cercano al origen es $\bar{x} = (1, 0, 0)^\top$. El Jacobiano de la restricción evaluado en $\bar{x}$ es $\nabla h(\bar{x}) = (-2, 0, 0)^\top$, por lo que se cumple LICQ. El subespacio tangente en $\bar{x}$ es el plano vertical $x_1 = 1$ generado por $d = (0, d_2, d_3)^\top$. El gradiente de la función objetivo en $\bar{x}$ es $\nabla f(\bar{x}) = (1, 0, 0)^\top$. El producto interno $\interno{\nabla f(\bar{x})}{d} = 0$. De este modo se verifica la condición necesaria: el gradiente es ortogonal al subespacio de direcciones tangentes.
\end{example}

\begin{corollary}[Singularidad Topológica por Exceso de Restricciones]\label{cor:punto_aislado}
    Supóngase que el punto $\bar{x} \in D$ satisface la condición LICQ y que la dimensión de las restricciones es mayor o igual que la dimensión del espacio primal ($l \ge n$). Entonces, $\bar{x}$ es un \textbf{punto aislado} del conjunto factible $D$, y es un minimizador local estricto para cualquier función continua $f$.
\end{corollary}

\begin{proof}
    Por el Teorema del Rango de Álgebra Lineal, $\dim(\ker J_h(\bar{x})) = n - \rk J_h(\bar{x})$. Bajo LICQ, el rango es máximo: $\rk J_h(\bar{x}) = \min(l, n) = n$. Sustituyendo, se tiene que $\dim(\ker J_h(\bar{x})) = 0$, lo que implica $\ker J_h(\bar{x}) = \{0\}$. Aplicando el Teorema del Subespacio Tangente (\cref{thm:subespacio_tangente}), los conos topológicos se reducen a $\mathcal{T}_D(\bar{x}) = \mathcal{B}_D(\bar{x}) = \{0\}$. Al no existir direcciones tangentes no nulas, no es posible definir una sucesión en $D \setminus \{\bar{x}\}$ que converja a $\bar{x}$. Por lo tanto, existe una vecindad $U$ de $\bar{x}$ tal que $D \cap U = \{\bar{x}\}$. Al ser el único punto factible en dicha vecindad, la condición $f(\bar{x}) \le f(x)$ se cumple trivialmente para todo $x \in D \cap U$.
\end{proof}